% !TEX program = lualatex
% Auto-generated from syllabus — 03: 情報を取り扱う上でのルール・マナー

\documentclass[aspectratio=43,professionalfonts,handout,10pt]{beamer}
\usepackage{beamer_template}

\newcommand{\LectureNum}{03}
\newcommand{\LectureTitle}{情報を取り扱う上でのルール・マナー}
\newcommand{\CourseName}{情報リテラシー}
\newcommand{\TextPages}{-}
\newcommand{\NextTopic}{課題の見つけ方と解決プロセス}
\newcommand{\NextPages}{pp.6-11}

\begin{document}
% --- Slide 1: title ---

\begin{frame}{\CourseName\quad 第\LectureNum 回「\LectureTitle 」}
  {\small \TermName\quad \TeacherName\quad ／\quad 教科書 \TextPages}

  \vfill

  \begin{block}{今日の流れ}
    \begin{enumerate}
      \item 情報モラルとは
      \item SNSの炎上とデマのリスク
      \item 個人情報の保護
      \item 著作権と知的財産
      \item 安全な情報発信のルール
    \end{enumerate}
  \end{block}
  \vfill

  \begin{block}{今日の目標}
    \begin{itemize}
      \item 情報や通信に関連する法令や規則等と，その必要性を説明できる。
      \item 情報社会で生活する上でのマナー，モラルの重要性について説明できる。
    \end{itemize}
  \end{block}
  \vfill

  \begin{exampleblock}{キーワード}
    \small \textbf{情報モラル} ／ \textbf{個人情報} ／ \textbf{著作権} ／ \textbf{肖像権} ／ \textbf{デジタルタトゥー}
  \end{exampleblock}
\end{frame}

% --- Slide 2: 情報モラルとは ---

\begin{frame}{情報モラルとは}
  \begin{block}{定義}
    \begin{itemize}
      \item 情報社会で適切に行動するための考え方やルール
      \item インターネット上での発言・投稿には責任が伴う
      \item 他人の権利（プライバシー・著作権・肖像権）を尊重する
    \end{itemize}
  \end{block}

  \vfill

  \begin{exampleblock}{なぜ情報モラルが必要か}
    \begin{itemize}
      \item 情報の\kw{複製性}：一度発信した情報はコピーされ拡散する
      \item 情報の\kw{残存性}：削除しても完全には消えない（デジタルタトゥー）
      \item 情報の\kw{伝播性}：瞬時に世界中に広がる可能性がある
    \end{itemize}
  \end{exampleblock}
\end{frame}

% --- Slide 3: SNSの炎上とデマ ---

\begin{frame}{SNSの炎上とデマ}
  \begin{block}{炎上の主な原因}
    \begin{itemize}
      \item 差別的・攻撃的な発言や投稿
      \item 根拠のないデマ情報の拡散
      \item 他人のプライバシーを侵害する投稿（無断で写真を公開など）
      \item 軽い気持ちの投稿が予想外に拡散される
    \end{itemize}
  \end{block}

  \vfill

  \begin{noticebox}[注意]
    \begin{itemize}
      \item 「あとで消せばOK」ではない → スクリーンショットで記録され残り続ける
      \item 匿名でも特定されるリスクがある（IPアドレス・投稿履歴の分析）
    \end{itemize}
  \end{noticebox}
\end{frame}

% --- Slide 4: 演習1 ---

\begin{frame}{演習1：この投稿，大丈夫？}
  \begin{block}{シナリオ}
    あなたは友達と一緒に撮った写真をSNSに投稿しようとしている。\\
    しかし，写真には知らない人も映っている。
  \end{block}

  \vfill

  \begin{alertblock}{考えてみよう}
    \begin{enumerate}
      \item この投稿にはどんなリスクがあるか？
      \item 投稿する前にどんな対策が必要か？
      \item もし投稿後にトラブルが起きたらどうするか？
    \end{enumerate}
  \end{alertblock}

  \vfill

  {\small 隣の人と2分間話し合い，意見をまとめよう。}
\end{frame}

% --- Slide 5: 個人情報の保護 ---

\begin{frame}{個人情報の保護}
  \begin{block}{個人情報とは}
    \begin{itemize}
      \item 氏名・住所・電話番号・メールアドレス・学籍番号など
      \item 特定の個人を識別できる情報すべてが該当する
      \item 顔写真・位置情報・購買履歴なども含まれる
    \end{itemize}
  \end{block}

  \vfill

  \begin{block}{個人情報保護法のポイント}
    \begin{itemize}
      \item 個人情報の取得には\kw{利用目的の明示}が必要
      \item 本人の同意なく\kw{第三者への提供}は原則禁止
      \item 漏えい時には本人への通知と個人情報保護委員会への報告が義務
    \end{itemize}
  \end{block}
\end{frame}

% --- Slide 6: 著作権と知的財産 ---

\begin{frame}{著作権と知的財産}
  \begin{block}{著作権の基本}
    \begin{itemize}
      \item 文章・音楽・写真・動画・プログラムなどの\kw{創作物}に自動的に発生する権利
      \item 著作権者の許可なく複製・配布・公開は原則\kw{違法}
      \item 保護期間：著作者の死後70年（日本の場合）
    \end{itemize}
  \end{block}

  \vfill

  \begin{exampleblock}{引用のルール}
    \begin{itemize}
      \item 引用部分を「」や引用ブロックで明確に区別する
      \item 出典（著者名・タイトル・URL等）を明記する
      \item 自分の文章が「主」，引用部分が「従」であること
    \end{itemize}
  \end{exampleblock}
\end{frame}

% --- Slide 7: 肖像権と関連する権利 ---

\begin{frame}{肖像権と関連する権利}
  \begin{block}{肖像権}
    \begin{itemize}
      \item 自分の顔や姿を無断で撮影・公開されない権利
      \item 法律に明文化されていないが，判例で認められている
      \item 友達の写真でもSNSに投稿する際は\kw{本人の許可}が必要
    \end{itemize}
  \end{block}

  \vfill

  \begin{block}{その他の関連する権利}
    \begin{itemize}
      \item \kw{プライバシー権}：私生活上の情報をみだりに公開されない権利
      \item \kw{パブリシティ権}：有名人の氏名・肖像の商業利用を管理する権利
    \end{itemize}
  \end{block}
\end{frame}

% --- Slide 8: 安全な情報発信のルール ---

\begin{frame}{安全な情報発信の5つのルール}
  \begin{block}{情報発信のポイント}
    \begin{enumerate}
      \item \kw{個人情報を安易に公開しない}（自分も他人も）
      \item \kw{投稿前に一度考える}：「この内容，公開しても大丈夫？」
      \item \kw{感情的な投稿を避ける}：怒りや悲しみのままの投稿は後悔のもと
      \item \kw{情報の正確性を確認する}：裏付けのない情報を拡散しない
      \item \kw{著作権・肖像権を尊重する}：他人の作品や写真は許可を得て使う
    \end{enumerate}
  \end{block}

  \vfill

  \begin{noticebox}[覚えておこう]
    匿名でも責任がある。拡散された情報は消せない。
  \end{noticebox}
\end{frame}

% --- Slide 9: 演習2 ---

\begin{frame}{演習2：情報発信のルールを考えよう}
  \begin{block}{グループワーク（5分議論 ＋ 2分まとめ）}
    \begin{enumerate}
      \item SNSを使う上で気をつけるべきことを3つ挙げよ
      \item 炎上やデマを防ぐためのポイントを考えよ
      \item グループで意見をまとめ，発表の準備をせよ
    \end{enumerate}
  \end{block}

  \vfill

  {\small 各グループ30秒で発表。}
\end{frame}

% --- Slide 10: 法令のまとめ ---

\begin{frame}{情報に関する主な法令}
  \begin{block}{知っておくべき法令}
    \begin{tabular}{ll}
      \hline
      \textbf{法令} & \textbf{主な内容} \\
      \hline
      個人情報保護法 & 個人情報の適切な取り扱い \\
      著作権法 & 創作物の権利保護 \\
      不正アクセス禁止法 & 他人のID・パスワードの不正使用を禁止 \\
      プロバイダ責任制限法 & 権利侵害時の情報開示請求 \\
      \hline
    \end{tabular}
  \end{block}

  \vfill

  \begin{exampleblock}{ポイント}
    「知らなかった」では済まされない。法令を理解し，正しく行動することが大切。
  \end{exampleblock}
\end{frame}

% --- チェックリスト ---

\begin{frame}{まとめ・チェックリスト}
  \begin{block}{自己チェック（できたら $\checkmark$ をつけよう）}
    \begin{itemize}
      \item[$\square$] 情報モラルの意味と必要性を説明できる
      \item[$\square$] SNSの炎上・デマのリスクと対策を理解した
      \item[$\square$] 個人情報保護法の基本を理解した
      \item[$\square$] 著作権・肖像権のルールを説明できる
      \item[$\square$] 安全な情報発信の5つのルールを理解した
    \end{itemize}
  \end{block}

  \vfill

  {\small 質問があれば授業後またはTeamsで受け付けます。}
\end{frame}

\end{document}
